% SOURCE_STATEMENT: "Establish purely global geometric inequalities for material bodies and determine whether their sharp forms require convexity, roundedness, or a related shape condition."
% Preparation and final self-audit: CONTRIBUTING.md prompts C and D; source/audit artifacts retained in the local physics_statements_2026_10_11 batch.
% !TeX program = lualatex
\documentclass[11pt]{article}
\usepackage[margin=25mm]{geometry}
\usepackage{fontspec}
\setmainfont{Latin Modern Roman}
\usepackage{amsmath,amssymb,mathtools}
\usepackage{unicode-math}
\setmathfont{Latin Modern Math}
\usepackage{xurl,hyperref}
\hypersetup{colorlinks=true,urlcolor=blue}
\setcounter{secnumdepth}{0}
\setlength{\parindent}{0pt}
\setlength{\parskip}{5pt}
\setlength{\emergencystretch}{3em}
\newcommand{\sref}[2]{\hyperlink{source:#1:#2}{[S#2]}}
\newcommand{\cataloguescope}[1]{\paragraph{Recorded target.}#1}
\newcommand{\R}{\mathbb R}
\newcommand{\C}{\mathbb C}
\newcommand{\N}{\mathbb N}
\newcommand{\Z}{\mathbb Z}
\newcommand{\Q}{\mathbb Q}
\newcommand{\F}{\mathbb F}
\newcommand{\eps}{\varepsilon}
\begin{document}
% UnsolvedMath ID 30003919; source catalogue code OWR-16411-012
\setcounter{section}{30003918}
\section{Global Geometric Inequalities for Material Bodies}\label{30003919}
{\small UnsolvedMath ID 30003919 · Mathematical Physics}
\subsection{Definitions and mathematical statement}

\paragraph{Shared notation.}$\mathbb N=\{1,2,\ldots\}$; the standard number fields have their usual meanings. All additional conventions are stated locally.


\paragraph{Definitions and conventions.}
A material body is a bounded matter region in gravitational initial data or a spacetime, with mass, angular momentum, charge and geometric size defined by a stated convention. A purely global inequality compares global quantities, unlike a bound confined to a single horizon section. Shape restrictions such as convexity or roundedness control how a radius or size describes the whole body. A sharp inequality includes the best constant and, where appropriate, an extremal or equality description.

\paragraph{Target retained from the canonical source catalogue.}
\paragraph{Catalogue wording retained for traceability.}
Establish purely global geometric inequalities for material bodies and determine whether their sharp forms require convexity, roundedness, or a related shape condition.

\paragraph{Source scope and qualification.}
No single radius, matter model or inequality is fixed by this programme. Such choices are part of the source question and cannot be filled by a generic Euclidean isoperimetric formula.
\subsection{Short English statement}
Find sharp global bounds for gravitating material bodies and decide which shape restrictions their validity requires.
\subsection{Sources}
\begin{itemize}
\item[\hypertarget{source:30003919:1}{S1}] ULAM.AI and the UnsolvedMath Contributors, UnsolvedMath, record 30003919 (OWR-16411-012). Catalogue curation under CC BY 4.0; definitions and English exposition added locally.
\url{https://huggingface.co/datasets/ulamai/UnsolvedMath}.
\item[\hypertarget{source:30003919:2}{S2}] Mathematical General Relativity (2018), Oberwolfach Reports, DOI 10.4171/owr/2018/36, pp. 2157--2251. Primary report accessed 2026-10-11; relevant PDF pages 79.
\url{https://doi.org/10.4171/owr/2018/36}.
\item[\hypertarget{source:30003919:3}{S3}] Geometrical inequalities review (catalogue-listed context; not independently consulted in this preparation).
\url{https://doi.org/10.1007/s41114-018-0014-7}.
\end{itemize}
\label{end:30003919}
\end{document}
